\section{The reservation discipline}

The discipline adds one object and one decision to the serving stack. The
object is the reservation: a ledger entry binding a flow to a guaranteed
resource vector for a lifetime. The decision is admission: whether a flow
enters the small reserved subset, enters an aggregate service class, or is
declined early. Everything else in this section is the precise semantics
of those two things, together with the supervisory and trust elements that
complete the contract. We emphasize at the outset what the discipline
keeps: the gateway accounting of Section 2 remains the identity and
metering substrate, and the scheduling engines of Section 2 remain the
execution machinery. The discipline sits at admission and constrains what
those layers may do to an admitted flow; it replaces neither.

A flow arrives with a descriptor: an estimated invocation count, a
requested cumulative output-token budget, an estimated cache footprint, a
service class, and a provenance class. Admission to the reserved subset
requires five conditions jointly, of which the load-bearing ones are that
the subset holds fewer than K flows, that projected commitment over a
forecast window stays under a configured ceiling, and that some device
set can hold the flow's cache footprint under the pinning rules. A flow
failing any condition is served in an aggregate class or declined early,
which is the cheapest moment in a flow's life to say no.

An admitted reservation covers three resources jointly, and the jointness
is the point, because a flow starved on any one of them fails entirely: a
compute share sufficient for its class latency budget, a cumulative
output-token budget metered at decode across every invocation of the
flow, and key-value cache pinned against eviction and migration for the
reservation lifetime, including across the idle gaps between invocations.
The budget is simultaneously a floor and a bound. Budgeted capacity is
committed, so the flow cannot be silently starved out of it, and the flow
cannot exceed it, so a runaway is a contained event rather than a
discovered one. Pinning has two embodiments, strict and tiered, the
latter spilling to a guaranteed restore path between invocations; the
evaluation prices both without favour.

The guarantee is stated negatively, and the reason is a claim about
auditability rather than style: for the reservation lifetime, no
resource-motivated action of the serving stack may preempt the flow's
invocations, evict or migrate its pinned cache, or reclaim its committed
budget. A prohibition can be audited for violations; a permission cannot
be audited for its non-occurrence. The sole override is a halt issued by
the management plane, which is a safety revocation under designated
authority and definitionally not scheduler preemption, since no capacity
condition can trigger it. A guarantee with a capacity-triggered escape
clause is a priority, and priorities are what the schedulers of Section 2
already have.

All other traffic runs in aggregate service classes with latency budgets
and scheduling weights, degraded under pressure by a probabilistic ladder
whose probability is keyed to committed headroom computed from the
reservation ledger rather than to observed queue depth. That keying is a
design claim the evaluation tests directly, not an assumption it rests
on.

Two elements complete the contract and were never economic claims. A
management plane holds pre-authorized halt, steer, and escalate
authority in a failure and trust domain separate from the inference
stack, which yields four properties that hold by construction: admission
is where organizational policy binds to a flow, the budget fixes any
runaway's worst case in advance rather than in an incident review, the
reservation ledger is the audit record, and the halt does not route
through the model being doubted. Two of the four are measured here
rather than asserted: Section 6 reports what the budget does under
adversarial injection, and what the difference between a priority and a
contract costs the flows subject to it. Trust-keyed admission gates
reservation eligibility on the provenance of the content driving a flow,
fixed before that content is interpreted. Appendix A carries the full
specification: the admission test, the ledger invariants, the
budget-exhaustion vocabulary, lease and multi-tenant semantics, and the
survivability embodiment.
